\section{Limitations and Conclusion}
\label{sec:conclusion}

\paragraph{Limitations.} All comparisons are single runs, and several change one setting on top of recipes that differ from each other. The downstream backbone predates the fixed-GSD stems, the dense context loss, the convolutional stems, the multi-observation encoder and the state head, so the recipe of \cref{sec:method} is not yet evaluated downstream. The cross-sample gap cannot separate content from location or product identity, and it compares only runs whose targets are alike. The mechanism of \cref{eq:decomp} is supported by probes and ablations but not proven. The corpus counts predate the seam repair, and polar footprints and the 1\,m NAC scale, where much of the scientific demand lies, are outside the runs reported here.

\paragraph{Conclusion.} After the raw products are calibrated and projected, the lunar archive becomes a dataset only once they are cut to common ground, and we described a pipeline that does this at the scale of the whole Moon while keeping every camera observation of a footprint, with its geometry, rather than one view per tile. On this corpus, predicting one instrument from another in latent space, under the sun of the target, is a workable pretraining task. Its frozen features detect craters about as well as a masked-token lunar model, gain substantially from a second product and respond to the sun they are given. Getting there taught us more about failure than about architecture. A cross-modal JEPA can reach a very low loss by encoding where a token is instead of what the ground looks like, and loss, within-sample rank and VICReg all read this as progress. A cross-sample comparison exposes it for one cosine per target, and we would log it in any JEPA whose targets are only partly predictable from the context.
